% =====================================================================
%  thmsmith Stage -1 proposal skeleton.
%  Written automatically by the thmsmith TS pipeline before Stage -1 runs.
%  Locked parameters: qid=stat\_sa\_cate\_pointwise, specialization=v1
%
%  HOW TO USE THIS FILE
%  --------------------
%  - The preamble (everything above the `% --- BEGIN CONTENT ---` marker)
%    and the `\section{...}` headers below are fixed by the pipeline.
%    Do NOT rewrite or rename them; the Step 3 reviewer subagent and the
%    downstream Stage 0 / Stage 1 prompts grep for these section titles.
%  - Each section contains exactly one `% TODO(stage-1 ...): ...`
%    placeholder. Replace each TODO with the content described for that
%    section in `CausalSmith/tools/src/discovery/prompts/stage_neg1_2_core.txt`
%    ("REQUIRED OUTPUT FORMAT (Stage -1 proposal .tex)").
%  - The `\paragraph{Locked parameters.}` block in the metadata block
%    must pin the chosen `(qid, specialization, p, assignment, target,
%    regression)` precisely. If the proposal maps to the Q1 pool-mode,
%    reproduce that block verbatim from
%    `CausalSmith/doc/panel_formula_question_generator_note_with_common_prompts.tex`
%    (Q2..Q6 templates have been retired). Otherwise (free-form
%    `(qid, specialization)` — the default), define each parameter
%    explicitly here (no implicit conventions). Stage 0 quotes this block;
%    do not rephrase it after locking.
%  - Removing a section header, renaming a macro, or rewriting the
%    preamble is a regression: the file must still match this skeleton
%    after you finish.
% =====================================================================

\documentclass[11pt]{article}

\usepackage{amsmath,amssymb,amsthm,mathtools}
\usepackage{enumitem}
\usepackage[colorlinks=true,linkcolor=blue,citecolor=blue,urlcolor=blue]{hyperref}
\usepackage{natbib}

\newcommand{\E}{\mathbb{E}}
\renewcommand{\Pr}{\mathbb{P}}
\newcommand{\one}{\mathbf{1}}
\newcommand{\transpose}{\mathsf{T}}
\newcommand{\calH}{\mathcal{H}}
\newcommand{\calF}{\mathcal{F}}
\newcommand{\calR}{\mathcal{R}}
\newcommand{\R}{\mathbb{R}}
\newcommand{\plim}{\operatorname*{plim}}

\theoremstyle{definition}
\newtheorem{definition}{Definition}
\newtheorem{assumption}{Assumption}
\newtheorem{example}{Example}

\theoremstyle{plain}
\newtheorem{theorem}{Theorem}
\newtheorem{conjecture}{Conjecture}
\newtheorem{proposition}{Proposition}
\newtheorem{lemma}{Lemma}
\newtheorem{corollary}{Corollary}

\theoremstyle{remark}
\newtheorem{remark}{Remark}

\title{thmsmith Stage -1 proposal: \texttt{stat\_sa\_cate\_pointwise} / \texttt{v1}%
  \\ \large\textit{Structure-agnostic pointwise CATE minimax rate with a cross-fit local DR learner and a Le Cam converse}}
\author{thmsmith Stage -1 proposer}
\date{\today}

\begin{document}
\maketitle

% --- BEGIN CONTENT ---  (everything above is fixed by the pipeline)

\paragraph{Locked parameters.}
Locked tuple:
\[
(\texttt{qid},\texttt{specialization},\mathcal{P}_n,\theta_0,\text{sampling model},\text{target object},\hat\theta_n)
\]
equals
\[
(\texttt{stat\_sa\_cate\_pointwise},\texttt{v1},
\mathcal{P}_n^{\mathrm{SA\text{-}CATE}}(x_0,h_n,r_{g,n},r_{m,n},b_n,\epsilon),
\tau_P(x_0),\text{iid unconfounded binary-treatment observations},
\rho_n,\hat\theta_n^{\mathrm{locDR}}).
\]
Here $W_i=(Y_i,D_i,X_i)$ are iid, $D_i\in\{0,1\}$, $x_0$ is fixed, $K_{h_n}$ is a nonnegative localizer, $p_n=E_P K_{h_n}(X)$, and $\ell_n(x)=K_{h_n}(x)/p_n$.
The causal estimand is $\tau_P(x_0)=E_P[Y(1)-Y(0)\mid X=x_0]$ under consistency, conditional unconfoundedness, and overlap $e_P(X)\in[\epsilon,1-\epsilon]$.
The model $\mathcal{P}_n^{\mathrm{SA\text{-}CATE}}$ is the class of laws satisfying Assumptions~\ref{ass:causal}--\ref{ass:local-rich}: local bias bounded by $b_n$, local effective sample size $n p_n\to\infty$, and black-box local nuisance rates $r_{g,n},r_{m,n}$ measured in the $\ell_n$-weighted local $L_2(P)$ norm.
The target object is the matched minimax pointwise absolute-error rate
\[
  \rho_n=b_n+(n p_n)^{-1/2}+\sqrt{r_{g,n}r_{m,n}}.
\]
The estimator $\hat\theta_n^{\mathrm{locDR}}$ is the $K$-fold cross-fit localized doubly robust learner in Section~\ref{sec:setup}.
%   (Q2..Q6 templates have been retired). Otherwise (free-form qid — the

\begin{abstract}
This proposal asks for the sharp structure-agnostic rate for estimating the conditional average treatment effect at a fixed covariate value $x_0$ when the only primitive nuisance information is a pair of black-box local mean-square rates.
Existing structure-agnostic lower bounds cover global treatment-effect functionals, while existing pointwise CATE minimax theory imposes smoothness and builds estimators around that structure.
The kernel claim is that a normalized cross-fit local DR learner attains $b_n+(n p_n)^{-1/2}+\sqrt{r_{g,n}r_{m,n}}$.
A matching Le Cam converse should follow from three local perturbations: a localization-bias pair, a local sampling pair, and a mixed outcome-propensity pair calibrated to the product rate.
The result would give pointwise DR-Learner practice a benchmark that separates localization bias, effective sample size, and nuisance-product error without assuming a forest, sparse-linear, or Holder model.
\end{abstract}

\section{Introduction}
Conditional average treatment effects are used when treatment effects vary with covariates and the analyst wants a value at a clinically or economically meaningful profile. The point value $\tau(x_0)$ is local and nonregular, so the usual average-effect intuition does not by itself say what black-box nuisance learning can buy.

\citet{KennedyBalakrishnanRobinsWasserman2024MinimaxCATE} give pointwise CATE minimax rates in a Holder-smooth model, and \citet{JinSyrgkanis2024StructureAgnosticDR,JinSyrgkanis2025GeneralLinearFunctionalLowerBounds} give structure-agnostic lower-bound theory for global or general linear functionals. The gap is the localized causal point functional where the rate must include the local effective sample size and the product of two black-box nuisance rates. The kernel is a Stat theorem: a matching achievability and lower-bound statement for $\tau(x_0)$ under a localized structure-agnostic model.

The immediate consumer is the DR-Learner CATE program in \citet{Kennedy2023DRCATE}. Its generic decomposition explains why a product nuisance term appears, but it does not identify the pointwise product term as unavoidable without smoothness. A reader using local DR pseudo-outcomes would report the three terms in $\rho_n$ as the benchmark and would not claim a faster pointwise rate without extra structure.

\begin{itemize}[leftmargin=*]
\item Theorem 1: The localized AIPW score identifies the smoothed CATE target and has second-order product drift under local weighted norms. Gap: \citet{ChernozhukovEtAl2018DML} and \citet{Kennedy2023DRCATE} contain the standard orthogonal-score cancellation, while this theorem fixes the normalization for the pointwise local target used by Conjectures 1--2. Consumer: local DR-Learner implementations that need the exact pseudo-outcome target before rate analysis.
\item Conjecture 1: The cross-fit local DR learner attains $\rho_n=b_n+(n p_n)^{-1/2}+\sqrt{r_{g,n}r_{m,n}}$. Gap: \citet{KennedyBalakrishnanRobinsWasserman2024MinimaxCATE} gives a Holder rate, while \citet{JinSyrgkanis2024StructureAgnosticDR} gives global structure-agnostic rates. Consumer: pointwise heterogeneous-effect reporting in medicine and policy targeting.
\item Conjecture 2: Three Le Cam local pairs lower-bound minimax risk by a constant multiple of $\rho_n$. Gap: \citet{JinSyrgkanis2025GeneralLinearFunctionalLowerBounds} gives the closest general lower-bound technology, but not the localized CATE point functional with $p_n$ and $b_n$ channels. Consumer: minimax benchmarking for local DR-Learner implementations.
\end{itemize}
% item 4 in `CausalSmith/tools/src/discovery/prompts/stage_neg1_2_core.txt` for full details.
%
%   Paragraph 1 --- Question and main claim.
%     The scientific or methodological question (1--2 sentences), then the headline theorem /
%     conjecture in plain English (1--2 sentences). Do not re-paraphrase the kernel; one clean
%     statement.
%
%   Paragraph 2 --- What is new vs. closest prior art, and where it lives.
%     Name 1--2 closest published comparators (bibkeys; with section/equation/page anchors for
%     flagship-aspirant claims --- N-thin-anchor applies), state the precise gap closed against
%     those comparators (1--2 sentences), then locate the proposal in the relevant area in one
%     sentence: exact identification, partial identification, panel / linear-projection methods,
%     causal machine learning, or network / interference causal inference. Do NOT default to
%     panel. The kernel must still be expressible as an ExactID, PartialID, or Panel theorem.
%
%   Paragraph 3 --- Consumer.
%     Name ONE concrete downstream consumer (1--2 sentences): a specific open problem this
%     unblocks, an applied setting whose analysis becomes tractable, an empirical finding that
%     gains a sharp theoretical anchor, or a published result that becomes a clean corollary.
%     Cite the consumer by bibkey or applied domain (generic "enables future work" / "of broad
%     interest" framings trigger N-no-implication). Close with one sentence answering "if a
%     reader takes only the conclusion away, what changes in their research practice?"
%
%   Paragraph 4 --- Results at a glance.
%     An itemized or enumerated list with exactly one entry per `Theorem N` / `Conjecture N` in
%     \S8. Each entry is one line:
%       `Theorem/Conjecture N: <one-line statement>. Gap: <closest paper(s) and why this
%        differs --- bibkeys>. Consumer: <named downstream beneficiary --- bibkey or applied
%        domain>.`
%     This block is the single home for per-result gap + consumer information. \S8 itself
%     carries math only --- no per-result "Why this is new" block (retired).

\section{Related work}
\citet{KennedyBalakrishnanRobinsWasserman2024MinimaxCATE} is the closest writing and CATE-rate comparator because it proves pointwise minimax rates under smoothness and frames CATE optimality as a minimax question.
\citet{Kennedy2023DRCATE} is the closest upper-bound estimator comparator because it studies a two-stage doubly robust CATE estimator with a model-free error decomposition and local polynomial double-residual specialization.
\citet{BonviniKennedyDukesBalakrishnan2024StructureAgnostic} and \citet{BalakrishnanKennedyWasserman2023StructureAgnostic} are adjacent minimax lineage because they study structure-agnostic or model-free treatment-effect estimation with nuisance-rate and robustness tradeoffs, while the present proposal fixes a nonregular point $x_0$ and separates $p_n$, $b_n$, and the local nuisance-product channel.
\citet{JinSyrgkanis2024StructureAgnosticDR} is the closest structure-agnostic treatment-effect comparator because it proves optimality of doubly robust estimators under black-box nuisance rates for ATE, ATT, and weighted ATE.
\citet{JinSyrgkanis2025GeneralLinearFunctionalLowerBounds} is the closest lower-bound comparator because it extends structure-agnostic lower-bound logic to general functional estimation, while the present proposal localizes that logic to a nonregular CATE point value with an effective local sample size.
\citet{ChernozhukovEtAl2018DML} supplies the cross-fitting and Neyman-orthogonal expansion that converts nuisance estimation into a product remainder.
\citet{ChernozhukovNeweySingh2022LocalRiesz} treats nonregular local linear functionals through regularized Riesz representers and motivates the $(n p_n)^{-1/2}$ local variance term.
\citet{ChernozhukovNeweyQuintasMartinezSyrgkanis2021RieszRegression} is adjacent because it estimates Riesz representers automatically; it would become relevant if the localizer or representer were learned rather than fixed.
\citet{SemenovaChernozhukov2021DMLCATE} gives simultaneous inference for CATE projections using orthogonal signals, but its target is a projection or best linear predictor rather than the unprojected point $\tau(x_0)$.
\citet{KallusMaoUehara2024LDML} provides a localized-DML template for target-dependent nuisances, but its main theory concerns efficient inference for local quantile effects rather than a matched black-box minimax CATE rate.
\citet{CuiTchetgenTchetgen2024SelectiveDR} studies selective machine learning of doubly robust functionals and marks the natural adaptive extension after the fixed-rate benchmark is settled.
In-repo, \texttt{CausalSmith/doc/API.md} records the Stat cluster as the home for estimation-theory kernels with matching converses, and the closest prior proposal signal is \texttt{eid\_dml\_orthogonality\_frontier}, whose reviewer asked for primitive empirical-process upper bounds and minimax lower-bound arguments rather than assumed frontier channels.

\paragraph{Exemplars.}
Writing exemplar: \citet{KennedyBalakrishnanRobinsWasserman2024MinimaxCATE}, DOI 10.1214/24-AOS2369; its arc from CATE motivation to unresolved optimality to a minimax formulation is mirrored in the introduction and setup.
Theorem 1: \citet{ChernozhukovEtAl2018DML} (estimation-inference, DOI 10.1111/ectj.12097) and \citet{ChernozhukovNeweySingh2022LocalRiesz} (estimation-inference, DOI 10.1093/ectj/utac002) supply the orthogonal-score and local-Riesz notation.
Conjecture 1: \citet{Kennedy2023DRCATE} (estimation-inference, arXiv:2004.14497) supplies the DR pseudo-outcome skeleton, with the setting changed to normalized point localization and black-box nuisance rates.
Conjecture 2: \citet{JinSyrgkanis2024StructureAgnosticDR} (estimation-inference, arXiv:2402.14264) and \citet{JinSyrgkanis2025GeneralLinearFunctionalLowerBounds} (estimation-inference, arXiv:2512.17341) supply the structure-agnostic lower-bound construction, and \citet{KennedyBalakrishnanRobinsWasserman2024MinimaxCATE} supplies the local CATE lower-bound geometry.
%   - Cite all named published papers necessary to understand the position of this work.
%   - Each citation MUST be accompanied by one sentence of relevance (no bibliography dumps).
%   - Cite relevant in-catalogue thmsmith proposals (by qid + headline statement) and `doc/API.md`
%     entries when they constitute in-repo prior art. Do NOT cite Lean files (file.lean:thm_name) ---
%     Lean substrate is not the novelty axis at this stage; Stage 1 owns formalization feasibility.
%   - This is the literature checklist the Step 3 reviewer will verify --- be accurate.

\section{Formal setup}\label{sec:setup}
Let $W=(Y,D,X)$ have law $P$, where $D\in\{0,1\}$ and $X$ takes values in a metric space $\mathcal{X}$.
Potential outcomes $Y(1),Y(0)$ satisfy consistency and conditional unconfoundedness.
Write
\[
  g_a(x)=E_P[Y\mid D=a,X=x],\qquad e(x)=P(D=1\mid X=x),
\]
and $\tau_P(x)=g_1(x)-g_0(x)$.
For a fixed point $x_0$ and a nonnegative kernel $K_h(x)=K(d(x,x_0)/h)$, set
\[
  p_h=E_P K_h(X),\qquad \ell_h(x)=K_h(x)/p_h,\qquad
  \theta_h(P)=E_P[\ell_h(X)\tau_P(X)].
\]
The pointwise target is $\theta_0(P)=\tau_P(x_0)$, with localization bias $b_h(P)=|\theta_h(P)-\theta_0(P)|$.

For nuisance functions $\eta=(g_0,g_1,e)$ define the localized AIPW score
\[
  \psi_h(W;\eta)=\ell_h(X)\left\{g_1(X)-g_0(X)
  +\frac{D(Y-g_1(X))}{e(X)}
  -\frac{(1-D)(Y-g_0(X))}{1-e(X)}\right\}.
\]
Given $K$ folds, train nuisance estimates $\hat\eta_{-k}=(\hat g_{0,-k},\hat g_{1,-k},\hat e_{-k})$ off fold $k$ and define
\[
  \hat\theta_n^{\mathrm{locDR}}
  =\frac{1}{n}\sum_{k=1}^{K}\sum_{i\in I_k}\psi_{h_n}(W_i;\hat\eta_{-k}).
\]
The local black-box nuisance rates are
\[
  R_{g,n}^2=E_P[\ell_{h_n}(X)\{(\hat g_1-g_1)^2+(\hat g_0-g_0)^2\}],\quad
  R_{m,n}^2=E_P[\ell_{h_n}(X)(\hat e-e)^2],
\]
with $R_{g,n}^2\le r_{g,n}$ and $R_{m,n}^2\le r_{m,n}$ in probability.
The target rate is $\rho_n=b_n+(n p_n)^{-1/2}+\sqrt{r_{g,n}r_{m,n}}$, where $p_n=p_{h_n}$ and $b_n$ bounds $b_{h_n}(P)$.

\begin{quote}\small
Locked tuple:
\[
(\texttt{qid},\texttt{specialization},\mathcal{P}_n,\theta_0,\text{sampling model},\text{target object},\hat\theta_n)
\]
equals
\[
(\texttt{stat\_sa\_cate\_pointwise},\texttt{v1},
\mathcal{P}_n^{\mathrm{SA\text{-}CATE}}(x_0,h_n,r_{g,n},r_{m,n},b_n,\epsilon),
\tau_P(x_0),\text{iid unconfounded binary-treatment observations},
\rho_n,\hat\theta_n^{\mathrm{locDR}}).
\]
Here $W_i=(Y_i,D_i,X_i)$ are iid, $D_i\in\{0,1\}$, $x_0$ is fixed, $K_{h_n}$ is a nonnegative localizer, $p_n=E_P K_{h_n}(X)$, and $\ell_n(x)=K_{h_n}(x)/p_n$.
The causal estimand is $\tau_P(x_0)=E_P[Y(1)-Y(0)\mid X=x_0]$ under consistency, conditional unconfoundedness, and overlap $e_P(X)\in[\epsilon,1-\epsilon]$.
The model $\mathcal{P}_n^{\mathrm{SA\text{-}CATE}}$ is the class of laws satisfying Assumptions~\ref{ass:causal}--\ref{ass:local-rich}: local bias bounded by $b_n$, local effective sample size $n p_n\to\infty$, and black-box local nuisance rates $r_{g,n},r_{m,n}$ measured in the $\ell_n$-weighted local $L_2(P)$ norm.
The target object is the matched minimax pointwise absolute-error rate
\[
  \rho_n=b_n+(n p_n)^{-1/2}+\sqrt{r_{g,n}r_{m,n}}.
\]
The estimator $\hat\theta_n^{\mathrm{locDR}}$ is the $K$-fold cross-fit localized doubly robust learner in Section~\ref{sec:setup}.
\end{quote}
%   - Define the data-generating environment for the chosen anchor cluster:
%       * Exact identification: SWIG / DAG / counterfactual setup, treatment + outcome + covariates, target estimand
%         (ATE / ATT / LATE / CDE / NDE / etc.), positivity / overlap / consistency assumptions as needed.
%       * Partial identification: observable distribution P, unobserved quantity \theta, identified set
%         \Theta_I(P), bounding functionals, and the restrictions (shape / monotonicity / IV / sensitivity) that
%         carve \Theta_I.
%       * Panel: panel index, treatment, finite-memory potential outcomes, response surface,
%         untreated baseline (E[Y_it(0_p) | D_i] = alpha_i + lambda_t unless locked parameters override).
%   - Reproduce locked parameters in a fenced environment exactly as in Stage 0.

\section{Assumptions}\label{sec:assumptions}
\begin{assumption}[Causal sampling and overlap]\label{ass:causal}
The observations $W_i=(Y_i,D_i,X_i)$ are iid. Consistency holds, $(Y(1),Y(0))\perp D\mid X$, and there is $\epsilon\in(0,1/2)$ such that $\epsilon\le e(X)\le 1-\epsilon$ almost surely. This is the unconfounded binary-treatment setup used by \citet{ChernozhukovEtAl2018DML} and \citet{Kennedy2023DRCATE}.
\end{assumption}

\begin{assumption}[Bounded local score]\label{ass:bounded}
There is $M<\infty$ such that $|Y|\le M$, $|g_a(X)|\le M$, $0\le K_{h_n}(X)\le 1$, $p_n=E K_{h_n}(X)>0$, and $E[\ell_{h_n}(X)^2]\asymp p_n^{-1}$. This mirrors the bounded-score and local-Riesz norm controls in \citet{ChernozhukovNeweySingh2022LocalRiesz}.
\end{assumption}

\begin{assumption}[Local bias modulus]\label{ass:bias}
For the chosen localization sequence $h_n$, $|\theta_{h_n}(P)-\tau_P(x_0)|\le b_n$ uniformly over $P\in\mathcal{P}_n^{\mathrm{SA\text{-}CATE}}$. This is the high-level local approximation condition used in pointwise minimax form by \citet{KennedyBalakrishnanRobinsWasserman2024MinimaxCATE}.
\end{assumption}

\begin{assumption}[Black-box nuisance rates]\label{ass:nuisance}
The cross-fit nuisance learners satisfy $R_{g,n}^2=O_P(r_{g,n})$ and $R_{m,n}^2=O_P(r_{m,n})$ in the local weighted norms of Section~\ref{sec:setup}, with $r_{g,n},r_{m,n}\to0$. This is the structure-agnostic oracle-rate convention of \citet{JinSyrgkanis2024StructureAgnosticDR} localized to $x_0$.
\end{assumption}

\begin{assumption}[Local effective sample size]\label{ass:effective}
$n p_n\to\infty$, $p_n\to0$ is allowed, and the cross-fit folds have sizes proportional to $n$. This is the triangular-array local-smoothing convention in \citet{KennedyBalakrishnanRobinsWasserman2024MinimaxCATE}.
\end{assumption}

\begin{assumption}[Local lower-bound richness]\label{ass:local-rich}
The model contains nondegenerate local neighborhoods in which bounded perturbations of the point CATE, the local mean mass, and the two nuisance functions can be made inside the support of $K_{h_n}$ while preserving Assumptions~\ref{ass:causal}--\ref{ass:effective}. This is the local analogue of the perturbability conditions used in \citet{JinSyrgkanis2025GeneralLinearFunctionalLowerBounds}.
\end{assumption}
%   - Each named assumption in its own `assumption` environment with a unique label.
%   - Every hypothesis cited in any Theorem/Conjecture must appear here.

\section{Main results}
\begin{theorem}[Theorem 1: localized AIPW identity]\label{thm:localized-aipw}
Under Assumptions~\ref{ass:causal} and~\ref{ass:bounded}, the localized AIPW score satisfies
\[
  E_P[\psi_h(W;\eta_P)]=\theta_h(P).
\]
For any nuisance tuple $\tilde\eta=(\tilde g_0,\tilde g_1,\tilde e)$ with $\tilde e\in[\epsilon/2,1-\epsilon/2]$, the standard product drift obeys
\[
  \left|E_P[\psi_h(W;\tilde\eta)]-\theta_h(P)\right|
  \le C_\epsilon
  \left(E_P[\ell_h(X)\{(\tilde g_1-g_1)^2+(\tilde g_0-g_0)^2\}]\right)^{1/2}
  \left(E_P[\ell_h(X)(\tilde e-e)^2]\right)^{1/2}.
\]
\end{theorem}
The calculation is the usual AIPW conditioning argument with the localization weight kept inside the conditional expectation, followed by Cauchy--Schwarz in the probability measure with density $\ell_h$ relative to $P$.

\begin{conjecture}[Conjecture 1: cross-fit local DR upper rate (NOT YET PROVED)]\label{conj:upper}
Under Assumptions~\ref{ass:causal}--\ref{ass:effective}, the estimator $\hat\theta_n^{\mathrm{locDR}}$ satisfies
\[
  \left|\hat\theta_n^{\mathrm{locDR}}-\tau_P(x_0)\right|
  =O_P\!\left(b_n+(n p_n)^{-1/2}+\sqrt{r_{g,n}r_{m,n}}\right)
\]
uniformly over $\mathcal{P}_n^{\mathrm{SA\text{-}CATE}}$.
\end{conjecture}
The statement commits to a concrete estimator and a single pointwise rate; regimes with $n p_n$ bounded are outside the locked model.

\begin{conjecture}[Conjecture 2: localized Le Cam converse (NOT YET PROVED)]\label{conj:lower}
Under Assumptions~\ref{ass:causal}--\ref{ass:local-rich}, there exists a constant $c>0$ such that every estimator $\tilde\theta_n$ obeys
\[
  \sup_{P\in\mathcal{P}_n^{\mathrm{SA\text{-}CATE}}}
  E_P\left|\tilde\theta_n-\tau_P(x_0)\right|
  \ge c\left(b_n+(n p_n)^{-1/2}+\sqrt{r_{g,n}r_{m,n}}\right).
\]
\end{conjecture}
The conjecture asserts the matching converse for the same locked model, using three local two-point channels rather than an estimator-specific obstruction.
%   - Number each result `Theorem N` (provable from existing substrate or by routine algebra)
%     or `Conjecture N` (novel kernel, verification still open --- mark `(NOT YET PROVED)`).
%   - State the math only --- assumptions, statement, and (for `Theorem`) a short proof sketch.
%     Surrounding prose <= 2 sentences per result. Do NOT attach a `\paragraph{Why this is new
%     and worth studying.}` block here; that block is retired. Per-result gap + consumer
%     information lives once in \S4 Paragraph 4 ("Results at a glance").
%   - At least one `Conjecture` is expected. All-`Theorem` usually means the proposer slipped statements past themselves.

\section{Sanity-check corollaries}
If $K_{h_n}\equiv1$, then $p_n=1$ and $b_n=0$, so Conjectures~\ref{conj:upper}--\ref{conj:lower} reduce to the global weighted-effect scale $n^{-1/2}+\sqrt{r_{g,n}r_{m,n}}$. This is the structure-agnostic WATE benchmark because the local weight has become a fixed bounded weight.

If the nuisance functions are known, then $r_{g,n}=r_{m,n}=0$ and the rate reduces to $b_n+(n p_n)^{-1/2}$. This is the ordinary local mean scale: deterministic localization bias plus variance from $n p_n$ effective observations.

If $\tau_P(x)$ is constant on the support of $K_{h_n}$ around $x_0$, then $b_n=0$. The result then says that pointwise CATE is no harder than a local weighted ATE except for the local sample-size loss and the nuisance product.
%   - At least one reduction the proposal must satisfy under existing assumptions
%     (homogeneous effects, no carry-over, T=2, iid, etc.). Catches "novelty by accident of error".
%   - FLAGSHIP-ONLY (SC10 exhibited-focal-object): For every flagship-aspirant headline
%     that names a constructive focal object (frontier T(P), threshold tau*, sharp set S*,
%     witness P*, bridge-function class H, identification algorithm, certificate,
%     comparator-divergence DGP), add an `Exhibit 9.k` block that COMPUTES the object on
%     a worked small instance (T=3 panel / 2-arm finite-support treatment / named small DAG)
%     from \S6 primitives --- NOT from \S7 assumed quantities. Include an internal-consistency
%     check (no over-determined system, no impossible cardinality, no zero denominator).
%     For non-equivalence claims between named estimators, the exhibit must produce
%     numerically (or symbolically) distinct estimator values on the instance and verify the
%     difference is not a normalization artifact. Missing exhibit fires `N-promissory-object`
%     at \S-0.5 and caps publishability at field.

\section{Proof-sketch route}
The upper route decomposes $\hat\theta_n^{\mathrm{locDR}}-\tau_P(x_0)$ into localization bias, a cross-fit empirical average of the oracle local score, and the conditional drift from replacing $\eta_P$ by $\hat\eta$. The oracle score term is controlled by Bernstein or Lindeberg inequalities at variance scale $E\ell_{h_n}^2/n\asymp(n p_n)^{-1}$. The drift term uses the supporting localized AIPW identity fold by fold and the black-box local nuisance rates, following the DML expansion in \citet{ChernozhukovEtAl2018DML} and the DR CATE pseudo-outcome route in \citet{Kennedy2023DRCATE}. The lower route adapts the structure-agnostic perturbation method of \citet{JinSyrgkanis2024StructureAgnosticDR,JinSyrgkanis2025GeneralLinearFunctionalLowerBounds}: construct one local-bias pair, one local mean perturbation of mass $p_n$, and one asymmetric outcome-propensity perturbation inside the $K_{h_n}$ neighborhood. Le Cam's two-point lemma converts bounded total variation into risk, and the three lower-bound channels are pasted with same-sign CATE shifts.
%   - One paragraph sketching the *mathematical* verification route: which classical
%     theorems / identities the proof would reduce to, and which boundary cases need
%     explicit algebra (e.g., a T=3 weight computation for panel kernels, a small-support
%     enumeration for partial-ID bounds, a back-door adjustment closed-form check for ExactID).
%     No proofs.
%   - FLAGSHIP-ONLY (SC11 definitional-unfold screen): For each flagship-aspirant headline,
%     write the reduction as an explicit chain:
%     `headline <- [unfold \S6 defs] <- [substitute \S7 assumptions] <- [step 1] <- ... <- [classical result]`.
%     Count k = the number of substantive non-trivial math steps that are NOT definition
%     unfolds and NOT direct \S7 assumption applications. Flagship requires k >= 3; otherwise
%     fire `C-definitional-unfold` and cap at field. Reflexive applications of standard
%     machinery (LP duality, FWL, support-function transform, do-calculus rule 2) count as
%     zero substantive steps; non-routine lemma applications count as one only if you explain
%     why a top referee would find them non-obvious.
%   - Do NOT name Mathlib lemmas, AutoID Lean files, or any Lean / formalization artefacts here.
%     Stage 1 (NL formalizer) owns the Lean-feasibility gate; the proposer's job is to convince
%     a *journal referee* the math is tractable, not a Lean engineer the substrate exists.

\section{Honest scope flags}
Theorem~\ref{thm:localized-aipw} is standard and serves only as a supporting prelude. Conjecture~\ref{conj:upper} still needs the empirical-process step for triangular-array local weights and cross-fit nuisance learners. Conjecture~\ref{conj:lower} is the main open content; Assumption~\ref{ass:local-rich} names the primitive richness needed for the perturbations but does not itself exhibit the three Le Cam pairs. The proposal is field-tier because it does not yet derive a universal lower-bound packing for every possible black-box oracle model.
%   - What is open vs proved.
%   - What was withdrawn during v1 --> vK iteration.
%   - Any assumptions the proposer is unsure about.

\section{Iteration trail}
\begin{itemize}[leftmargin=*]
\item v1 (cold-start) -- initial draft of the field-tier Stat kernel; prior verdict: none.
\item v2 (revise) -- restored the localized AIPW identity as Theorem 1, aligned the Section 4 result inventory, and corrected the selective-DR Biometrika citation; prior verdict: REVISE.
\end{itemize}
%   - Bulleted summary of v1, v2, ..., vK with reviewer verdict / one-line change at each step.
%   - Even on a v1 ACCEPT, include a one-line entry for transparency.

\section*{Bibliography}
\begin{thebibliography}{99}
\bibitem[Balakrishnan, Kennedy, and Wasserman(2023)]{BalakrishnanKennedyWasserman2023StructureAgnostic}
Balakrishnan, S., Kennedy, E. H., and Wasserman, L. (2023).
Model-free and structure-agnostic treatment-effect estimation.
arXiv preprint.

\bibitem[Cui and Tchetgen Tchetgen(2024)]{CuiTchetgenTchetgen2024SelectiveDR}
Cui, Y. and Tchetgen Tchetgen, E. J. (2024).
Selective machine learning of doubly robust functionals.
\emph{Biometrika}, 111(2), 517--535.

\bibitem[Bonvini, Kennedy, Dukes, and Balakrishnan(2024)]{BonviniKennedyDukesBalakrishnan2024StructureAgnostic}
Bonvini, M., Kennedy, E. H., Dukes, O., and Balakrishnan, S. (2024).
Structure-agnostic minimax analysis for treatment-effect estimation with flexible nuisance learning.
arXiv preprint.

\bibitem[Chernozhukov et~al.(2018)]{ChernozhukovEtAl2018DML}
Chernozhukov, V., Chetverikov, D., Demirer, M., Duflo, E., Hansen, C., Newey, W., and Robins, J. (2018).
Double/debiased machine learning for treatment and structural parameters.
\emph{The Econometrics Journal}, 21(1), C1--C68.

\bibitem[Chernozhukov, Newey, and Singh(2022)]{ChernozhukovNeweySingh2022LocalRiesz}
Chernozhukov, V., Newey, W. K., and Singh, R. (2022).
Debiased machine learning of global and local parameters using regularized Riesz representers.
\emph{The Econometrics Journal}, 25(3), 576--601.

\bibitem[Chernozhukov et~al.(2021)]{ChernozhukovNeweyQuintasMartinezSyrgkanis2021RieszRegression}
Chernozhukov, V., Newey, W. K., Quintas-Martinez, V., and Syrgkanis, V. (2021).
Automatic debiased machine learning via Riesz regression.
arXiv:2104.14737.

\bibitem[Jin and Syrgkanis(2024)]{JinSyrgkanis2024StructureAgnosticDR}
Jin, J. and Syrgkanis, V. (2024).
Structure-agnostic optimality of doubly robust learning for treatment effect estimation.
arXiv:2402.14264.

\bibitem[Jin and Syrgkanis(2025)]{JinSyrgkanis2025GeneralLinearFunctionalLowerBounds}
Jin, J. and Syrgkanis, V. (2025).
Sharp structure-agnostic lower bounds for general functional estimation.
arXiv:2512.17341.

\bibitem[Kallus, Mao, and Uehara(2024)]{KallusMaoUehara2024LDML}
Kallus, N., Mao, X., and Uehara, M. (2024).
Localized debiased machine learning: Efficient inference on quantile treatment effects and beyond.
\emph{Journal of Machine Learning Research}, 25(16), 1--59.

\bibitem[Kennedy(2023)]{Kennedy2023DRCATE}
Kennedy, E. H. (2023).
Towards optimal doubly robust estimation of heterogeneous causal effects.
\emph{Electronic Journal of Statistics}, 17(2), 3008--3049.

\bibitem[Kennedy et~al.(2024)]{KennedyBalakrishnanRobinsWasserman2024MinimaxCATE}
Kennedy, E. H., Balakrishnan, S., Robins, J. M., and Wasserman, L. (2024).
Minimax rates for heterogeneous causal effect estimation.
\emph{The Annals of Statistics}, 52(2), 793--816.

\bibitem[Semenova and Chernozhukov(2021)]{SemenovaChernozhukov2021DMLCATE}
Semenova, V. and Chernozhukov, V. (2021).
Debiased machine learning of conditional average treatment effects and other causal functions.
\emph{The Econometrics Journal}, 24(2), 264--289.
\end{thebibliography}

\end{document}
